\chapter{Introduction}
\label{sec:introduction}

Invasive alien plants alter ecosystems, reduce native biodiversity and
are costly to control once they have spread over large areas
\cite{ipbes2023}. \emph{Acacia dealbata} (silver wattle), native to
south-eastern Australia \cite{dessi2021}, is one such species in
Portugal, where the whole genus \emph{Acacia} is listed as invasive
under national law \cite{dl92-2019}. It forms stands of its own and
spreads inside stands of other species \cite{icnf2019}, and knowing where it currently grows is a prerequisite for any control
effort. Producing an up-to-date map over a large area is nonetheless
difficult: the species must be told apart from visually similar trees
and shrubs in a satellite image, the land it occupies changes from
year to year, and the reference data needed to train and check a
classifier are slow and expensive to collect.

This work, carried out within the MAP.INVADER project on Earth
Observation methods for mapping invasive alien plants, investigates
whether two recent developments, citizen-science observations and
Earth Observation foundation models, can reduce the dependency on that
expert-collected reference data for mapping \emph{A.~dealbata} in the
Coimbra district. This intermediate report reviews the relevant
literature, defines the approach and its evaluation, and reports a
first experiment.

\section{Context and Motivation}
\label{sec:introduction-motivation}

Knowing where the species grows is precisely what the official
sources do not provide. Portugal's national forest inventory, the
reference source for forest composition, is produced once a decade
from photo-points and field plots; its latest edition has 2015 as
reference year and was published in 2019 \cite{icnf2019}. It recorded
8.4 thousand hectares of stands dominated by acacias, three times the
1995 figure, but as a sample it locates none of them, and it counts
the species only where it dominates, although it also occurred in
14\% of the maritime-pine area and 19\% of the eucalyptus area. For a
species that spreads inside other stands and re-establishes after
fire, a decennial sample is not a map.

Keeping a map current matters more for a species that benefits from
disturbance. In burnt stands in the Coimbra district, fire promoted
\emph{A.~dealbata} seed germination and no treatment prevented
regeneration within three years \cite{riveiro2025}, so a burnt area
where the species was present is a likely site of reinvasion rather
than a cleared one. Large fires are recurrent: in 2024 Portugal was
the most affected country in the European Union, with 147 thousand
hectares burnt, almost all of them in one week of September and most
of them in the centre of the country \cite{sanmiguel2025}. A map of
current stands that is cheap to produce and to repeat is what makes it
possible to anticipate where control will be needed, instead of
recording the invasion after it has spread.

Satellites make it possible to observe vegetation repeatedly and at
scale. The Sentinel-2 mission, for example, revisits the same location
every few days and records several spectral bands beyond what the
human eye sees, at a resolution of 10 metres per pixel. Mapping a
species from these images means classifying each pixel as that species
or not, using whatever signal best separates it from its surroundings;
for \emph{A.\ dealbata}, this is mainly its distinctive winter
flowering period \cite{domingo2023,martins2016}.

Classifying pixels requires reference data: locations where the
species is known to be present (and, ideally, known to be absent),
used to train and evaluate the classifier. Reference data of this kind
are usually produced by specialists: field visits to confirm the
species, or careful interpretation of very-high-resolution images by
someone who knows what to look for. The result is reliable and
delimits stands rather than single points, but it is slow and
expensive to collect, it is produced once, for one area and one year,
and every new area or new year costs the same expert time again. This
is what keeps species-level maps rare and quickly outdated. Citizen
science offers the opposite trade-off: on platforms such as
iNaturalist, members of the public record species observations with a
photograph, a date and a location of limited precision, identified by
other users rather than by an expert \cite{gillespie2024}. These
records are opportunistic rather than designed, but they are collected
continuously and over the whole country at no cost to the mapping
effort. If these observations could do part of the work that expert
data does today, maps could be produced for more areas and updated
more often than expert surveys allow.

What makes this plausible now is the arrival of \emph{Earth
Observation foundation models}: models pre-trained on large volumes of
satellite imagery that turn each pixel into a fixed-length vector, or
\emph{embedding}, capturing structure learned across many places and
times, rather than a hand-engineered spectral feature. TESSERA
\cite{feng2025} and AlphaEarth \cite{brown2025} are two recent
examples. For a reader familiar with machine learning but not Earth
Observation, the closest analogy is a pre-trained language or vision
model: instead of engineering features by hand, a light classifier is
trained on top of embeddings produced by a model that has already
learned general-purpose structure from far more data than any single
study could label.

The two developments fit each other. Citizen science supplies labels
cheaply, but few of them and of uneven quality; foundation-model
embeddings let a light classifier learn from few labels, because the
general structure of the imagery has already been learned without
them \cite{ball2026}. Whether that combination actually works for a
specific invasive species in a specific region is an open empirical
question, not something that follows automatically from either
technology on its own.

Using citizen-science observations as reference data raises problems
of its own. Observations cluster along roads and trails rather than
where the species grows \cite{dimson2023}; in the Coimbra district, a
single observer made 90\% of the 2\,465 presences kept for this work,
so the data may cover few of the conditions in which
\emph{A.~dealbata} occurs. They record presences only, never a
confirmed absence \cite{phillips2009}. Their positions are imprecise
relative to a 10-metre pixel, and many were made in years other than
the imagery \cite{ball2026}. And because nearby observations resemble
each other, a model evaluated on points next to its training points
looks better than it is \cite{karasiak2022}, while classifiers trained
on foundation-model embeddings lose accuracy when applied to a region
other than the one they were trained in \cite{yang2026}.

\section{Research Questions and Goals}
\label{sec:introduction-rqs}

Each ingredient above -- foundation models for mapping, citizen
science as reference data, and an invasive plant as the target -- has
been studied on its own or in pairs, but their combination, with an
explicit measure of how far citizen-science data can stand in for
expert data, has not, as the literature reviewed in the State of the
Art (Chapter~\ref{sec:state-of-the-art}) shows. The main
question of this dissertation is:

\begin{quote}
To what extent can citizen-science observations, turned into
reference data and represented by Earth Observation foundation-model
embeddings, reduce the amount of expert reference data needed to
train a model that maps \emph{Acacia dealbata} in the Coimbra
district?
\end{quote}

The reduction is read from a curve: the performance of a model trained
on citizen-science observations alone and with increasing fractions of
the expert polygons added, compared with a model trained on expert
data alone, all measured on the same held-out locations, which hold
presences from both sources. The answer is a degree, not a yes/no.
Answering this question is the general objective of this
dissertation. Four secondary questions operationalise it, in order of
priority, each followed by the specific objective that answers it.
Each uses only what the questions before it define, and the one that
depends on Sentinel-2 imagery comes last.

\begin{itemize}
  \item \textbf{RQ1 -- Citizen science as a substitute for expert
  data.} To what extent does a model trained only on citizen-science
  observations reach the performance of models trained with expert
  reference data?\\
  \emph{Objective:} measure the performance gained by adding 10\% to
  100\% of the expert polygons to a citizen-science base, compare it
  with an expert-only model, and report the performance of the
  citizen-science-only model on expert reference data.
  \item \textbf{RQ2 -- From point observations to reference data.} To
  what extent do the grade and accuracy filters applied to
  citizen-science observations, and the extension of each label from
  the observed pixel to the vegetation patch around it, change the
  performance of \emph{A.~dealbata} mapping?\\
  \emph{Objective:} build and document the process that turns
  citizen-science point observations into reference data, and measure
  how its grade and accuracy filters, and the extension of each label
  to the surrounding vegetation patch, change performance.
  \item \textbf{RQ3 -- Transferability away from the sampled area.}
  To what extent does the performance of the citizen-science-trained
  model drop from locations near the sampled area to locations far
  from it?\\
  \emph{Objective:} measure the difference in performance between
  locations near and far from the sampled area.
  \item \textbf{RQ4 -- Foundation models against conventional
  features.} To what extent do TESSERA embeddings improve
  \emph{A.~dealbata} mapping over Sentinel-2 features when trained on
  cleaned citizen-science data, with all observations and with few?\\
  \emph{Objective:} compare TESSERA-based and Sentinel-2-based models,
  with all observations and with few.
\end{itemize}

Figure~\ref{fig:introduction-pipeline} places each question at the
step of the pipeline it acts on. Each citizen-science observation
becomes a labelled pixel, the pixel is represented by a
foundation-model embedding, a presence-only classifier is trained on
those embeddings, and the classifier produces the species map. RQ2
acts on the passage from observation to label, RQ4 on the
representation, and RQ3 on the map. Expert polygons stay outside the
main flow: they are used as test data alongside citizen science and,
for RQ1, are added to the training data in increasing fractions.

\begin{figure}[htb]
\centering
\resizebox{\textwidth}{!}{%
\begin{tikzpicture}[
  node distance=4mm and 8mm,
  stage/.style={draw, rounded corners=2pt, align=center, minimum height=10mm,
    text width=25mm, font=\footnotesize, inner sep=2pt},
  expert/.style={stage, dashed},
  note/.style={align=center, font=\scriptsize, text width=25mm},
  rq/.style={font=\scriptsize\bfseries},
  flow/.style={-{Latex[length=2mm]}, thick},
  dflow/.style={flow, dashed}
]
\node[stage] (obs) {Citizen-science\\observation};
\node[stage, right=of obs] (lab) {Reference\\pixel label};
\node[stage, right=of lab] (emb) {Embedding\\(TESSERA)};
\node[stage, right=of emb] (clf) {Presence-only\\classifier};
\node[stage, right=of clf] (map) {\emph{A.~dealbata} map\\+ confidence};
\node[expert, below=9mm of clf] (exp) {Expert polygons};

\node[note, below=1mm of obs] {iNaturalist,\\research grade};
\node[note, below=1mm of lab] {10~m pixel, 2023};
\node[note, below=1mm of emb] {EO foundation\\model};
\node[note, below=1mm of clf, xshift=-4mm, text width=16mm] {presence vs.\\background};
\node[note, below=1mm of map] {Coimbra district};
\node[note, below=1mm of exp] {test data;\\added to training\\0--100\%};

\draw[flow] (obs) -- (lab);
\path (obs.north east) -- (lab.north west) node[rq, midway, above=1mm] {RQ2};
\draw[flow] (lab) -- (emb);
\draw[flow] (emb) -- (clf);
\draw[flow] (clf) -- (map);
\draw[dflow] ([xshift=6mm]exp.north) -- node[rq, right=2pt] {RQ1} ([xshift=6mm]clf.south);
\node[rq, above=1mm of emb] {RQ4};
\node[rq, above=1mm of map] {RQ3};
\end{tikzpicture}}
\caption{Overview of the proposed pipeline. Each research question is
placed at the step it acts on; RQ4 replaces the embedding step by
Sentinel-2 features.}
\label{fig:introduction-pipeline}
\end{figure}

Two choices narrow the thesis proposal, which named a limited number
of species and the fine-tuning of foundation models. The work maps a
single species, \emph{A.~dealbata}, because the expert polygons cover
only that species, so the comparison with expert data is possible only
for it. And it uses the foundation model frozen, as a source of
embeddings, because the TESSERA evaluation found no benefit from
fine-tuning its encoder (Section~\ref{sec:background-fm-embeddings}).

RQ1 alone constitutes a complete result and is the core deliverable
of this work, answering the main objective of the thesis proposal,
developed within the MAP.INVADER project; a preliminary version of its
curve is reported in the preliminary experiment of the Proposed Work
(Section~\ref{sec:proposed-work-preliminary}). The
full design of each question, including hypotheses, validation and
metric, is given in the experimental design of the Proposed Work
(Chapter~\ref{sec:proposed-work}). Besides these
answers, the work produces the \emph{A.~dealbata} map of the study
area, with its confidence map, as a product of the second semester:
it is presented as what the model predicts, not as a validated map,
since no ground truth covers the whole area.

Table~\ref{tab:introduction-problems} returns to the problems raised in
the motivation (Section~\ref{sec:introduction-motivation}) and says where each one is
addressed.

\begin{table}[htb]
\centering
\begin{tabularx}{\textwidth}{|>{\raggedright\arraybackslash}X|>{\raggedright\arraybackslash}X|>{\raggedright\arraybackslash}X|}
\hline
\textbf{Problem} & \textbf{In machine-learning terms} & \textbf{Addressed in} \\
\hline
Observations cluster near roads and trails \cite{dimson2023} &
Sampling / dataset bias & RQ3 \\
\hline
A single observer made most observations (83\% overall, 90\% of the
cleaned presences used in this dissertation) & Low training-data diversity &
RQ3; observer balancing (Proposed Work, Chapter~\ref{sec:proposed-work}) \\
\hline
Citizen science records only presences, never confirmed absences
\cite{phillips2009} & Positive-unlabelled learning & Uniform-random
background pool (Proposed Work, Chapter~\ref{sec:proposed-work}) \\
\hline
Imprecise locations and mixed pixels \cite{ball2026} & Label noise &
RQ2 \\
\hline
Limited volume of reference data & Data efficiency & RQ1; RQ4 \\
\hline
Observations recorded in years other than the satellite imagery
\cite{ball2026} & Temporal mismatch between data and imagery & Single
reference year; disturbance filter (Proposed Work, Chapter~\ref{sec:proposed-work}) \\
\hline
Nearby observations inflate accuracy estimates \cite{karasiak2022} &
Data leakage / spatial autocorrelation & Spatial validation
(Proposed Work, Chapter~\ref{sec:proposed-work}) \\
\hline
Weaker transfer of embeddings to new regions \cite{yang2026} &
Distribution shift & RQ3 \\
\hline
\end{tabularx}
\caption{Practical problems in using citizen-science data with Earth
Observation foundation models, and where each is addressed in this
dissertation.}
\label{tab:introduction-problems}
\end{table}

\section{Document Structure}
\label{sec:introduction-structure}
This document reports the work of the first semester and is
organised in five chapters.

\begin{itemize}
  \item \textbf{Chapter~\ref{sec:introduction}, Introduction:}
  presents the problem of mapping invasive plants in Portugal, the
  motivation for combining citizen-science observations with Earth
  Observation foundation models, the research questions and the goals
  of this dissertation.
  \item \textbf{Chapter~\ref{sec:background}, Background:} introduces
  the concepts the rest of the document relies on: \emph{Acacia
  dealbata} as an invasive plant in Portugal, Sentinel-2 imagery,
  Earth Observation foundation models and TESSERA, citizen-science
  data from iNaturalist, and presence-only learning with spatial
  validation.
  \item \textbf{Chapter~\ref{sec:state-of-the-art}, State of the
  Art:} reviews the literature on mapping invasive plants from
  satellite imagery, on citizen-science observations as reference
  data, and on foundation-model embeddings for species mapping, and
  identifies the gap this dissertation addresses.
  \item \textbf{Chapter~\ref{sec:proposed-work}, Proposed Work:}
  describes the proposed approach, the pipeline that turns
  citizen-science observations into reference labels, the
  experimental design for each research question, the datasets
  prepared during the first semester, and a preliminary result for
  RQ1. It ends with the work plan and the risk assessment.
  \item \textbf{Chapter~\ref{sec:conclusion}, Conclusion and Future
  Work:} summarises what was achieved in the first semester, what the
  preliminary result does and does not yet show, and the steps planned
  for the second semester.
\end{itemize}
